\justifying
\NSFCBodyText

\subsubsection{已具备的实验条件}

新疆大学计算机科学与技术学科具有多语言多模态认知计算与内容安全、信号检测与智能处理、大数据分析与智能控制等建设方向。可依托的科研平台包括教育部丝路多语言认知计算国际合作联合实验室、鹏城实验室新疆网络节点、新疆信号检测与处理重点实验室、新疆多语种信息技术重点实验室、国家语委新疆多语种信息技术研究中心及新疆多模态智能处理与信息安全工程技术研究中心，可为多模态、时序、网络与安全相关研究提供学科环境。

计算资源方面，学院人工智能算力中心材料记载拥有 GPU/NPU 训练及推理服务器/工作站 130 余台、各型加速卡 600 余块，FP32/FP16 精度下人工智能算力分别为 43/120 PFLOPS；2025 年购买 A100 8 卡服务器三年租赁服务。上述条件可以支撑模型训练、对照实验、弱网与资源策略仿真和数据处理。A100 服务的确切起止日期、申请人可用配额及能否覆盖本项目全周期，仍需在正式申报前核验，故不将其写成项目全周期专用算力。

\subsubsection{尚缺少的实验条件}

申请人确认目前可落实的是基础可控实验场景，弱网链路需要采用模拟方式构造。尚未形成可写入申请书的新疆口岸/中亚运输车辆/边境仓储合作准入、现场数据规模、统一货物标识下的运输—入库—仓储配对记录，以及振动、倾斜、电子封志、光感、温湿度、烟雾、门禁、视频、边缘网关和功耗设备的具体型号与数量。视频及跨境数据若进入后续试验，须先落实授权、脱敏和合规边界。计算平台不能替代上述现场与端侧条件。

\subsubsection{拟解决的途径}

项目按“基础可控异常—参数化弱网模拟—代表性边缘验证—条件成熟后的有限场景核验”分层落实。在可控环境中依据三时钟和证据质量规范完成异常触发、完整原始流和日志采集；利用参数化网络仿真或流量整形改变带宽、时延、丢包率与断连时长，开展 2G/3G 弱网机制试验，并将结论标注为模拟结果；在后续可获得的代表性板卡上测试时延、内存和通信字节，具备功耗仪后再报告实测能耗，否则只给出可复核的估算口径。有限现场及跨阶段验证仅在合作、授权、实体配对和统计功效满足后开展，拟合作和模拟续航结果不作为已落实条件或真实设备测量结果。

\subsubsection{利用研究基地的计划与落实情况}

依托多模态智能处理、信号检测、网络节点和人工智能算力条件，开展多源时序建模、视觉异常分析、网络/系统仿真和模型训练；边缘原型和物流场景试验按照设备与合作条件分阶段接入。项目将把训练数据、模型版本、弱网参数、资源日志和评估脚本统一归档，保证成对重放、分组测试和结果追溯；平台共享算力按校内实际审批与配额使用，不写成项目专用资源。
